\documentclass[10pt,a4paper]{article}
\usepackage{fontspec}
\setmainfont{Liberation Serif}

\usepackage[margin=0.6in]{geometry}
\usepackage{xcolor}
\usepackage{graphicx}
\usepackage{array}
\usepackage{enumitem}
\usepackage{booktabs}
\usepackage{caption}
\setmonofont{Inconsolatazi4}[Scale=0.9]

\setlength{\parindent}{0pt}
\setlength{\parskip}{6pt}

\begin{document}

% ======================== HEADER ========================
\begin{center}
{\large\textbf{TRIBHUVAN UNIVERSITY}}\\
Institute of Science and Technology\\
\vspace{0.2cm}
{\Large\textbf{SOFTWARE PROJECT MANAGEMENT}}\\
{\large\textbf{LABORATORY REPORT}}\\
Bachelor of Science in Computer Science and Information Technology (BSc. CSIT)\\
Seventh Semester\\
\vspace{0.2cm}
\textbf{Topic:} Project Scheduling, WBS/RBS, Gantt Chart and PERT/CPM Network Diagram using ProjectLibre
\end{center}

\vspace{0.2cm}
\noindent\textbf{Submitted by:}\\
Name: Samir Paudel\hfill Roll No./Symbol No: 114/079 --- 79010103\\
Program: BSc. CSIT, 7th Semester\\[0.1cm]
\textbf{Submitted to:}\\
Department of Computer Science and Information Technology

\hrulefill

% ======================== OBJECTIVE ========================
\section*{Objective}

The objective of this lab is to understand the fundamentals of project scheduling and network modeling used in software project management. Specifically, this lab covers:
\begin{itemize}[nosep]
    \item Creating a project schedule using project management software (MS Project / OpenProj / ProjectLibre).
    \item Constructing a Work Breakdown Structure (WBS) and Resource Breakdown Structure (RBS).
    \item Generating basic and relationship-based Gantt charts.
    \item Producing Task Usage and Resource Usage reports.
    \item Drawing PERT (Program Evaluation and Review Technique) and CPM (Critical Path Method) network diagrams and identifying the critical path.
\end{itemize}

% ============================================================
%                     TASK 1
% ============================================================
\newpage
\section*{Task 1: Student Profile Management System --- Project Scheduling}

Software used: ProjectLibre (open-source alternative to Microsoft Project). The same steps apply to MS Project / OpenProj with minor UI differences.

\subsection*{(a) Creating the Project}

\begin{enumerate}[nosep]
    \item Open ProjectLibre and select \texttt{File $\to$ New}.
    \item In the New Project dialog, enter the following details:
    \begin{itemize}[nosep]
        \item Project Name: \textbf{Student Profile Management}
        \item Start Date: \textbf{20/02/2024}
        \item Manager: \textbf{Samir Paudel}
        \item Organization: \textbf{Patan Multiple Campus}
    \end{itemize}
    \item Click \texttt{OK} to create a blank project shell.
    \item Set the project calendar to \texttt{Standard} from \texttt{Project $\to$ Project Information}, and confirm the working days/hours (8 hrs/day).
\end{enumerate}

[Insert Screenshot: New Project dialog box showing Project Name = \textit{Student Profile Management}]

\subsection*{(b) Creating the Work Breakdown Structure (WBS)}

\begin{enumerate}[nosep]
    \setcounter{enumi}{4}
    \item Switch to the Gantt Chart view (\texttt{View $\to$ Task Gantt}).
    \item Enter each activity name in the Task Name column, in outline order 1 to 6.
    \item Enter the Duration for every task in days.
    \item Enter the Start Date for Task 1 (20/02/2024); ProjectLibre auto-computes subsequent dates once predecessors are linked.
    \item In the Predecessors column, enter the task numbers as given (e.g. Task 4 $\to$ \texttt{2,3}).
    \item Double-click each task $\to$ Resources tab $\to$ assign the resource and the \% Units of involvement (e.g. Designer at 15\%).
\end{enumerate}

[Insert Screenshot: Task sheet / Gantt table filled with all 6 activities, durations, dates, predecessors and resources]

\vspace{0.2cm}
\noindent\textbf{WBS Data Table (as entered in ProjectLibre)}

\begin{table}[h!]
\centering
\renewcommand{\arraystretch}{1.15}
\begin{tabular}{|c|l|c|c|c|c|l|}
\hline
\textbf{S.N.} & \textbf{Activity} & \textbf{Duration (days)} & \textbf{Start Date} & \textbf{Finish Date} & \textbf{Predecessor} & \textbf{Resource (\% Involvement)} \\ \hline
1 & Planning \& Analysis & 7 & 2024-02-20 & 2024-02-26 & --- & Team Leader (100\%) \\ \hline
2 & Requirement Collection \& Definition & 15 & 2024-02-27 & 2024-03-12 & 1 & System Analyst (20\%) \\ \hline
3 & Designing & 7 & 2024-03-13 & 2024-03-19 & 2 & Designer (15\%) \\ \hline
4 & Development & 21 & 2024-03-20 & 2024-04-09 & 2, 3 & Programmer (70\%) \\ \hline
5 & Testing & 8 & 2024-04-10 & 2024-04-17 & 2, 3, 4 & Test Engineer (20\%) \\ \hline
6 & Deployment & 7 & 2024-04-18 & 2024-04-24 & 3, 4, 5 & --- \\ \hline
\end{tabular}
\caption*{WBS Data Table}
\end{table}

\noindent\textit{Note:} Dates are computed on a Finish-to-Start (FS) dependency basis --- a successor task begins the day after the latest finishing predecessor. Total project duration = 65 calendar days (20 Feb 2024 -- 24 Apr 2024).

\subsection*{(c) Project Scheduling --- Gantt Diagrams}

\begin{enumerate}[nosep]
    \setcounter{enumi}{10}
    \item With all tasks, durations and dependencies entered, ProjectLibre automatically draws bars in the Gantt Chart pane on the right.
    \item For the Basic Gantt Chart, view only task bars against the timeline (\texttt{View $\to$ Gantt Chart}).
    \item For the Relationship-based Gantt Chart, ensure predecessor links are entered so that dependency arrows (Finish-to-Start) are drawn connecting related task bars.
    \item Right-click the Gantt area $\to$ Bar Styles to customize colors (optional).
\end{enumerate}

\noindent\textbf{(i) Basic Gantt Chart}

Fig 1: Basic Gantt Chart --- Student Profile Management System

\vspace{0.3cm}
\noindent\textbf{(ii) Relationship-based Gantt Chart (with dependency arrows)}

Fig 2: Relationship-based Gantt Chart showing Finish-to-Start dependency arrows

\subsection*{(d) Work Breakdown Structure (WBS) and Resource Breakdown Structure (RBS) Diagrams}

\begin{enumerate}[nosep]
    \setcounter{enumi}{14}
    \item In ProjectLibre, go to \texttt{View $\to$ WBS Chart} (or, if unavailable, construct the WBS chart manually using the outline numbers assigned to each task, as shown below).
    \item The WBS decomposes the project into phases (Level 1) and further into deliverable-based work packages (Level 2).
    \item For the RBS, go to \texttt{View $\to$ Resource Sheet} to list all resources, then group them by category (Human Resources / Material Resources) to draw the RBS hierarchy.
\end{enumerate}

\noindent\textbf{Work Breakdown Structure (WBS)}

Fig 3a: WBS --- Phases 1 to 3

Fig 3b: WBS --- Phases 4 to 6

\noindent\textbf{Resource Breakdown Structure (RBS)}

Fig 4: Resource Breakdown Structure (RBS) --- Student Profile Management System

\subsection*{(e) Task Usage and Resource Usage Reports}

\begin{enumerate}[nosep]
    \setcounter{enumi}{17}
    \item Go to \texttt{View $\to$ Task Usage} to see each task with its assigned resource(s) and the time-phased work (hours).
    \item Go to \texttt{View $\to$ Resource Usage} to see each resource with the list of tasks it is assigned to and total work.
    \item Work is computed as: Work = Duration $\times$ Resource Units (e.g., Programmer at 70\% for 21 days $\times$ 8 hrs/day = 117.6 hrs).
\end{enumerate}

\vspace{0.2cm}
\noindent\textbf{Task Usage Report}

\begin{table}[h!]
\centering
\renewcommand{\arraystretch}{1.15}
\begin{tabular}{|c|l|c|c|c|c|}
\hline
\textbf{S.N.} & \textbf{Task} & \textbf{Duration} & \textbf{Resource} & \textbf{Units} & \textbf{Work (hrs)} \\ \hline
1 & Planning \& Analysis & 7d & Team Leader & 100\% & 56.0 \\ \hline
2 & Requirement Collection \& Definition & 15d & System Analyst & 20\% & 24.0 \\ \hline
3 & Designing & 7d & Designer & 15\% & 8.4 \\ \hline
4 & Development & 21d & Programmer & 70\% & 117.6 \\ \hline
5 & Testing & 8d & Test Engineer & 20\% & 12.8 \\ \hline
6 & Deployment & 7d & Not assigned & --- & 0.0 \\ \hline
\end{tabular}
\caption*{Task Usage Report}
\end{table}

\vspace{0.2cm}
\noindent\textbf{Resource Usage Report}

\begin{table}[h!]
\centering
\renewcommand{\arraystretch}{1.15}
\begin{tabular}{|c|l|c|c|}
\hline
\textbf{S.N.} & \textbf{Resource} & \textbf{Assigned Task(s)} & \textbf{Total Work (hrs)} \\ \hline
1 & Team Leader & Planning \& Analysis & 56.0 \\ \hline
2 & System Analyst & Requirement Collection \& Definition & 24.0 \\ \hline
3 & Designer & Designing & 8.4 \\ \hline
4 & Programmer & Development & 117.6 \\ \hline
5 & Test Engineer & Testing & 12.8 \\ \hline
\end{tabular}
\caption*{Resource Usage Report}
\end{table}

[Insert Screenshot: ProjectLibre Task Usage view and Resource Usage view]

\subsection*{(f) Network Modeling Diagram --- PERT Chart and CPM}

\begin{enumerate}[nosep]
    \setcounter{enumi}{20}
    \item Go to \texttt{View $\to$ Network Diagram (PERT Chart)} in ProjectLibre to view the Activity-on-Node (AON) network automatically generated from the task list and predecessors.
    \item Each node box shows the task ID, name, duration, and (once the CPM calculation is run via \texttt{Tools $\to$ Options} or the built-in scheduler) the Early Start (ES), Early Finish (EF), Late Start (LS) and Late Finish (LF).
    \item Critical tasks (zero slack/float) are usually highlighted in red by the tool.
    \item The Critical Path is the longest path through the network and determines the minimum project duration.
\end{enumerate}

Fig 5: PERT/CPM Network Diagram (AON) --- Student Profile Management System

\vspace{0.2cm}
\noindent\textit{Analysis:} Since every activity in this project depends directly on the completion of all its predecessors in a single sequential chain ($1\to2\to3\to4\to5\to6$), there is only one path through the network. Hence all six activities are critical, with zero slack, and the Critical Path is:

\vspace{0.2cm}
\noindent\textbf{Critical Path:} 1 $\to$ 2 $\to$ 3 $\to$ 4 $\to$ 5 $\to$ 6

\noindent\textbf{Project Duration = 7 + 15 + 7 + 21 + 8 + 7 = 65 days}

% ============================================================
%                     TASK 2
% ============================================================
\newpage
\section*{Task 2: PERT Diagram and Critical Path}

Given the following project activities:

\begin{table}[h!]
\centering
\renewcommand{\arraystretch}{1.15}
\begin{tabular}{|c|c|c|}
\hline
\textbf{Activity} & \textbf{Predecessor} & \textbf{Duration} \\ \hline
A & --- & 2 \\ \hline
B & A & 4 \\ \hline
C & A & 5 \\ \hline
D & B & 6 \\ \hline
E & C & 3 \\ \hline
F & D, E & 4 \\ \hline
\end{tabular}
\caption*{Project Activities}
\end{table}

\subsection*{Forward Pass (Earliest Start / Earliest Finish)}

\begin{itemize}[nosep]
    \item ES(A) = 0, EF(A) = 0 + 2 = 2
    \item ES(B) = EF(A) = 2, EF(B) = 2 + 4 = 6
    \item ES(C) = EF(A) = 2, EF(C) = 2 + 5 = 7
    \item ES(D) = EF(B) = 6, EF(D) = 6 + 6 = 12
    \item ES(E) = EF(C) = 7, EF(E) = 7 + 3 = 10
    \item ES(F) = max(EF(D), EF(E)) = max(12, 10) = 12, EF(F) = 12 + 4 = 16
\end{itemize}

\subsection*{Backward Pass (Latest Start / Latest Finish)}

\begin{itemize}[nosep]
    \item LF(F) = 16, LS(F) = 16 $-$ 4 = 12
    \item LF(D) = LS(F) = 12, LS(D) = 12 $-$ 6 = 6
    \item LF(E) = LS(F) = 12, LS(E) = 12 $-$ 3 = 9
    \item LF(B) = LS(D) = 6, LS(B) = 6 $-$ 4 = 2
    \item LF(C) = LS(E) = 9, LS(C) = 9 $-$ 5 = 4
    \item LF(A) = min(LS(B), LS(C)) = min(2, 4) = 2, LS(A) = 2 $-$ 2 = 0
\end{itemize}

\subsection*{Slack / Float Table}

\begin{table}[h!]
\centering
\renewcommand{\arraystretch}{1.15}
\begin{tabular}{|c|c|c|c|c|c|c|}
\hline
\textbf{Activity} & \textbf{ES} & \textbf{EF} & \textbf{LS} & \textbf{LF} & \textbf{Slack (LS$-$ES)} & \textbf{Critical?} \\ \hline
A & 0 & 2 & 0 & 2 & 0 & Yes \\ \hline
B & 2 & 6 & 2 & 6 & 0 & Yes \\ \hline
C & 2 & 7 & 4 & 9 & 2 & No \\ \hline
D & 6 & 12 & 6 & 12 & 0 & Yes \\ \hline
E & 7 & 10 & 9 & 12 & 2 & No \\ \hline
F & 12 & 16 & 12 & 16 & 0 & Yes \\ \hline
\end{tabular}
\caption*{Slack / Float Table}
\end{table}

Fig 6: PERT Network Diagram --- Activities A to F

\subsection*{Path Comparison}

\begin{itemize}[nosep]
    \item Path A--B--D--F = 2 + 4 + 6 + 4 = 16 days
    \item Path A--C--E--F = 2 + 5 + 3 + 4 = 14 days
\end{itemize}

\noindent\textbf{Critical Path:} A $\to$ B $\to$ D $\to$ F

\noindent\textbf{Project Duration = 16 days}

% ============================================================
%                     TASK 3
% ============================================================
\newpage
\section*{Task 3: AOA Network Diagram and Critical Path Calculation}

Given activities represented in (i, j) node notation:

\begin{table}[h!]
\centering
\renewcommand{\arraystretch}{1.15}
\begin{tabular}{|c|c|}
\hline
\textbf{Activity (i, j)} & \textbf{Duration (Days)} \\ \hline
(1, 2) & 4 \\ \hline
(1, 3) & 6 \\ \hline
(2, 4) & 5 \\ \hline
(3, 4) & 3 \\ \hline
(4, 5) & 2 \\ \hline
(3, 5) & 8 \\ \hline
\end{tabular}
\caption*{Project Activities}
\end{table}

\subsection*{Forward Pass --- Earliest Event Time E(j)}

\begin{itemize}[nosep]
    \item E(1) = 0
    \item E(2) = E(1) + 4 = 0 + 4 = 4
    \item E(3) = E(1) + 6 = 0 + 6 = 6
    \item E(4) = max[E(2)+5, E(3)+3] = max[4+5, 6+3] = max[9, 9] = 9
    \item E(5) = max[E(4)+2, E(3)+8] = max[9+2, 6+8] = max[11, 14] = 14
\end{itemize}

\subsection*{Backward Pass --- Latest Event Time L(i)}

\begin{itemize}[nosep]
    \item L(5) = 14
    \item L(4) = L(5) $-$ 2 = 14 $-$ 2 = 12
    \item L(3) = min[L(4)$-$3, L(5)$-$8] = min[12$-$3, 14$-$8] = min[9, 6] = 6
    \item L(2) = L(4) $-$ 5 = 12 $-$ 5 = 7
    \item L(1) = min[L(2)$-$4, L(3)$-$6] = min[7$-$4, 6$-$6] = min[3, 0] = 0
\end{itemize}

\subsection*{Event Time Table}

\begin{table}[h!]
\centering
\renewcommand{\arraystretch}{1.15}
\begin{tabular}{|c|c|c|c|c|}
\hline
\textbf{Event (Node)} & \textbf{Earliest Time E} & \textbf{Latest Time L} & \textbf{Slack (L$-$E)} & \textbf{Critical?} \\ \hline
1 & 0 & 0 & 0 & Yes \\ \hline
2 & 4 & 7 & 3 & No \\ \hline
3 & 6 & 6 & 0 & Yes \\ \hline
4 & 9 & 12 & 3 & No \\ \hline
5 & 14 & 14 & 0 & Yes \\ \hline
\end{tabular}
\caption*{Event Time Table}
\end{table}

Fig 7: AOA Network Diagram --- Nodes 1 to 5

\subsection*{Critical Path}

An activity (i, j) lies on the critical path only if E(i)=L(i), E(j)=L(j), and E(i)+duration=E(j). Applying this test:

\begin{itemize}[nosep]
    \item (1,3): E(1)+6 = 0+6 = 6 = E(3) and both nodes critical $\to$ \textbf{Critical}
    \item (3,5): E(3)+8 = 6+8 = 14 = E(5) and both nodes critical $\to$ \textbf{Critical}
    \item (1,2), (2,4), (3,4), (4,5): fail the equality test $\to$ Not critical
\end{itemize}

\noindent\textbf{Critical Path:} 1 $\to$ 3 $\to$ 5 (via activities (1,3) and (3,5))

\noindent\textbf{Project Duration = 6 + 8 = 14 days}

% ============================================================
%                     CONCLUSION
% ============================================================
\newpage
\section*{Conclusion}

This lab provided hands-on experience in software project scheduling using a project management tool (ProjectLibre) and manual network analysis. Key learnings include:

\begin{itemize}[nosep]
    \item Creating a project plan with a Work Breakdown Structure (WBS) linking tasks, durations, dates, dependencies and resources.
    \item Visualizing the schedule through basic and relationship (dependency-linked) Gantt charts.
    \item Decomposing the project into WBS and RBS hierarchies for clearer planning and resource allocation.
    \item Generating Task Usage and Resource Usage reports to track effort distribution.
    \item Applying the forward pass and backward pass technique of PERT/CPM to compute ES, EF, LS, LF, slack, and to identify the critical path --- the sequence of activities with zero slack that determines the minimum project completion time.
\end{itemize}

For the three exercises:
\begin{enumerate}[nosep]
    \item The \textbf{Student Profile Management} project has a critical path of \textbf{65 days} along all six sequential activities.
    \item The \textbf{PERT problem} has a critical path \textbf{A--B--D--F} of \textbf{16 days}.
    \item The \textbf{AOA network} has a critical path \textbf{1--3--5} of \textbf{14 days}.
\end{enumerate}

\vfill

\begin{flushright}
\textbf{Samir Paudel}\\
Roll No: 114/079\\
Section: D
\end{flushright}

\end{document}
